% RESUMO - PT
\begin{resumo}
  \vspace{-15pt}
  
  Este trabalho apresenta o desenvolvimento de uma aplicação web destinada ao apoio da gestão operacional de uma lanchonete. A proposta parte da identificação de dificuldades associadas ao uso de registros manuais, como erros no lançamento de pedidos, demora no atendimento, inconsistências no controle de produtos e limitação do acesso a informações gerenciais. O sistema foi estruturado com o framework Laravel e a linguagem PHP, utilizando MySQL para a persistência dos dados e Redis como mecanismo de cache. A arquitetura adota o padrão Model--View--Controller e o controle de acesso baseado em funções, com o objetivo de separar responsabilidades e restringir as funcionalidades conforme o perfil do usuário. Entre as funções previstas estão o cadastro e a autenticação de usuários, o gerenciamento de produtos, a apresentação do catálogo, o carrinho de compras e o acompanhamento dos pedidos. O uso de Docker contribui para a padronização do ambiente de execução, enquanto Git e GitHub apoiam o controle de versões. Espera-se que a solução reduza a ocorrência de falhas operacionais, centralize as informações do estabelecimento e proporcione maior agilidade ao atendimento e ao acompanhamento das atividades.

  Palavras-chave: \palavraschaveemlinha
\end{resumo}


% RESUMO - EN
\begin{resumo}[Abstract]
  \vspace{-15pt}
  
  \begin{otherlanguage*}{english}
  This work presents the development of a web application designed to support the operational management of a snack bar. The proposal is based on difficulties associated with manual records, including order-entry errors, service delays, inconsistencies in product control, and limited access to management information. The system was structured with the Laravel framework and PHP, using MySQL for data persistence and Redis as a caching mechanism. Its architecture follows the Model--View--Controller pattern and role-based access control to separate responsibilities and restrict features according to each user profile. The planned features include user registration and authentication, product management, catalog display, a shopping cart, and order tracking. Docker contributes to a standardized execution environment, while Git and GitHub support version control. The solution is expected to reduce operational failures, centralize business information, and improve the speed of customer service and activity monitoring.
  
  Keywords: \inlinekeywords
\end{otherlanguage*}
\end{resumo}
